% Содержимое данного документа позаимсвовано из Приложения Е из документа https://library.bsuir.by/m/12_101945_1_141950.pdf

% Содержимое данного документа позаимсвовано из Приложения Е из документа http://www.bsuir.by/m/12_113415_1_66883.pdf

\thispagestyle{empty}

\begin{singlespace}

{\small
  \begin{center}
    \begin{minipage}{0.8\textwidth}
      \begin{center}
        {\normalsize ОТЗЫВ}\\[1em]
        на дипломный проект\\студента факультета информационных технологий 
        и управления\\Учреждения образования <<Белорусский государственный университет информатики и радиоэлектроники>>\\
        %\fio \
        Веркович Елизаветы Васильевны\\
        на тему: <<\topicName>>
      \end{center}
    \end{minipage}
  \end{center}

Тема дипломного проекта, посвященного разработке персонального ассистента организации времени обучающихся, является актуальной, поскольку связана с задачами повышения эффективности самостоятельной учебной деятельности, снижения проявлений академической прокрастинации и внедрения интеллектуальных цифровых средств поддержки планирования в образовательный процесс.

В ходе дипломного проектирования Веркович~Е.\,В. в целом решила поставленную задачу: выполнен анализ предметной области и существующих подходов, разработаны требования к программному средству, спроектированы архитектура и модель данных, реализован программный комплекс, обеспечивающий декомпозицию учебных задач, формирование черновика расписания и его последующую корректировку пользователем, а также выполнено технико-экономическое обоснование разработки и использования системы.

В процессе работы студентка проявила достаточную степень самостоятельности и инициативности. Основные проектные и технические решения принимались ею осознанно, с аргументацией выбранных подходов, инструментов и архитектурных решений. Поставленные этапы выполнялись последовательно и в соответствии с календарным графиком.

Веркович~Е.\,В. продемонстрировала умение пользоваться специальной литературой, нормативными и методическими источниками, а также современными публикациями по тематике интеллектуальных ассистентов, организации времени, проектирования программных систем и применения языковых моделей. Материалы источников были не только изучены, но и корректно использованы при обосновании принятых решений.

В ходе выполнения дипломного проекта студентка проявила способности как к инженерной, так и к исследовательской работе. Это выразилось в умении формализовать задачу, сопоставить альтернативные варианты реализации, спроектировать структуру программного средства, реализовать взаимодействие серверных компонентов и пользовательского интерфейса, а также оценить практическую применимость и экономическую эффективность разработанного решения.

Пояснительная записка и графический материал оформлены аккуратно, структура работы логична, изложение материала последовательное. Дипломный проект соответствует заданию и отражает достаточный уровень подготовки автора по специальности.

Считаю, что Веркович~Е.\,В. подготовлена к самостоятельной профессиональной деятельности по специальности 1-40~03~01 <<Искусственный интеллект>> и может быть представлена к присвоению квалификации <<Инженер-системотехник>>.

  \vfill
  \noindent
  \begin{minipage}{0.54\textwidth}
    \begin{flushleft}
      Руководитель дипломного проекта:\\
      ассистент кафедры ИИТ\\
      \mbox{}\\
      \underline{\hspace*{2em}} \underline{\hspace*{6.5em}} \the\year{}~г.
    \end{flushleft}
  \end{minipage}
  \begin{minipage}{0.44\textwidth}
    \begin{flushright}
      \underline{\hspace*{3cm}} М.\,Г.~Соколович
    \end{flushright}
  \end{minipage}
 }

\end{singlespace}


\clearpage